\begin{table}[htbp]
\centering
\caption{The classical bounds test on the 1999-2007 regime.}
\label{tab:tab04_classical}
\begin{threeparttable}
\begin{tabular}{lrrrrrrrr}
\toprule
specification & n & k & F & t & I(0) 5\% & I(1) 5\% & verdict & theta(neer) \\
\midrule
dollar alone & 104 & 1 & 1.629000 & -1.128000 & 4.996000 & 5.973000 & fail to reject & -0.875000 \\
+ 7 controls fixed & 104 & 1 & 11.084000 & -4.697000 & 4.996000 & 5.973000 & reject & 0.004000 \\
\bottomrule
\end{tabular}

\begin{tablenotes}[flushleft]\small
\item Bounds are simulated at this sample size, not read off a table calibrated at T = 1000. Holding controls fixed is the best the classical procedure can do with them; it cannot account for their trend content.
\end{tablenotes}
\end{threeparttable}
\end{table}